\documentclass[UTF8,11pt]{ctexart}
\usepackage[a4paper,margin=24mm]{geometry}
\usepackage{amsmath,amssymb,amsthm,mathtools,mathrsfs}
\usepackage[all]{xy}
\usepackage{xurl,hyperref,enumitem}
\usepackage{tikz}
\usetikzlibrary{arrows.meta,cd,decorations.markings,calc,patterns}
\hypersetup{hidelinks,breaklinks=true}
\setlength{\emergencystretch}{3em}
\allowdisplaybreaks
\newtheorem*{theorem}{Theorem}
\newtheorem*{thm}{Theorem}
\newtheorem*{lemma}{Lemma}
\newtheorem*{proposition}{Proposition}
\newtheorem*{prop}{Proposition}
\newtheorem*{corollary}{Corollary}
\newtheorem*{cor}{Corollary}
\newtheorem*{claim}{Claim}
\theoremstyle{definition}
\newtheorem*{definition}{Definition}
\newtheorem*{defn}{Definition}
\newtheorem*{remark}{Remark}
\newtheorem*{example}{Example}
\newcommand{\R}{\mathbb R}
\newcommand{\C}{\mathbb C}
\newcommand{\Z}{\mathbb Z}
\newcommand{\Q}{\mathbb Q}
\newcommand{\N}{\mathbb N}
\newcommand{\F}{\mathbb F}
\newcommand{\D}{\mathbb D}
\newcommand{\bB}{\mathbb B}
\newcommand{\bC}{\mathbb C}
\newcommand{\bR}{\mathbb R}
\newcommand{\bZ}{\mathbb Z}
\newcommand{\bN}{\mathbb N}
\newcommand{\bQ}{\mathbb Q}
\newcommand{\bD}{\mathbb D}
\newcommand{\bF}{\mathbb F}
\newcommand{\CP}{\mathbb{CP}}
\newcommand{\RP}{\mathbb{RP}}
\newcommand{\sO}{\mathscr O}
\newcommand{\sI}{\mathscr I}
\newcommand{\sF}{\mathscr F}
\newcommand{\sG}{\mathscr G}
\newcommand{\sH}{\mathscr H}
\newcommand{\sR}{\mathscr R}
\newcommand{\sM}{\mathscr M}
\newcommand{\sV}{\mathscr V}
\newcommand{\sU}{\mathscr U}
\DeclareMathOperator{\Hom}{Hom}
\DeclareMathOperator{\Ext}{Ext}
\DeclareMathOperator{\Tor}{Tor}
\DeclareMathOperator{\Spec}{Spec}
\DeclareMathOperator{\Proj}{Proj}
\DeclareMathOperator{\supp}{supp}
\DeclareMathOperator{\im}{im}
\DeclareMathOperator{\id}{id}
\DeclareMathOperator{\Tr}{Tr}
\DeclareMathOperator{\tr}{tr}
\DeclareMathOperator{\rank}{rank}
\DeclareMathOperator{\ord}{ord}
\DeclareMathOperator{\dist}{dist}
\DeclareMathOperator{\End}{End}
\DeclareMathOperator{\Aut}{Aut}
\DeclareMathOperator{\Gal}{Gal}
\DeclareMathOperator{\vol}{vol}
\DeclareMathOperator{\Cl}{Cl}
\DeclareMathOperator{\coker}{coker}
\DeclareMathOperator{\codim}{codim}
\DeclareMathOperator{\divergence}{div}
\DeclareMathOperator{\Ric}{Ric}
\DeclareMathOperator{\grad}{grad}
\DeclareMathOperator{\Hess}{Hess}
\newcommand{\abs}[1]{\left\lvert#1\right\rvert}
\newcommand{\sabs}[1]{\lvert#1\rvert}
\newcommand{\babs}[1]{\bigl\lvert#1\bigr\rvert}
\newcommand{\Babs}[1]{\Bigl\lvert#1\Bigr\rvert}
\newcommand{\bbabs}[1]{\biggl\lvert#1\biggr\rvert}
\newcommand{\BBabs}[1]{\Biggl\lvert#1\Biggr\rvert}
\newcommand{\norm}[1]{\left\lVert#1\right\rVert}
\newcommand{\snorm}[1]{\lVert#1\rVert}
\newcommand{\bnorm}[1]{\bigl\lVert#1\bigr\rVert}
\newcommand{\Bnorm}[1]{\Bigl\lVert#1\Bigr\rVert}
\newcommand{\bbnorm}[1]{\biggl\lVert#1\biggr\rVert}
\newcommand{\BBnorm}[1]{\Biggl\lVert#1\Biggr\rVert}
\newcommand{\widebar}[1]{\overline{#1}}
\newcommand{\nicefrac}[2]{\frac{#1}{#2}}
\newcommand{\linnprod}[2]{\langle#1,#2\rangle}
\newcommand{\myindex}[1]{#1}
\newcommand{\glsadd}[1]{}
\newcommand{\avoidbreak}{}
\newcommand{\sourceref}[1]{\textnormal{[原文：\nolinkurl{#1}]}}
\newcommand{\thmref}[1]{\sourceref{#1}}
\newcommand{\lemmaref}[1]{\sourceref{#1}}
\newcommand{\propref}[1]{\sourceref{#1}}
\newcommand{\corref}[1]{\sourceref{#1}}
\newcommand{\defnref}[1]{\sourceref{#1}}
\newcommand{\exampleref}[1]{\sourceref{#1}}
\newcommand{\exerciseref}[1]{\sourceref{#1}}
\newcommand{\figureref}[1]{\sourceref{#1}}
\newcommand{\sectionref}[1]{\sourceref{#1}}
\newcommand{\appendixref}[1]{\sourceref{#1}}
\newcommand{\stacksref}[2]{\href{https://stacks.math.columbia.edu/tag/#1}{#2 [#1]}}


\newcommand{\E}{\mathbb E}
\newcommand{\Prob}{\mathbb P}
\newcommand{\Reff}{\mathcal R_{\mathrm{eff}}}
\newcommand{\eps}{\varepsilon}
\DeclareMathOperator{\diam}{diam}
\DeclareMathOperator{\Var}{Var}
\DeclareMathOperator{\Crit}{Crit}
\newcommand{\dd}{\,\mathrm d}

\begin{document}
\section*{073\quad 首个非零同伦群与同调群的Hurewicz同构}
\subsection*{来源与计数目标}
Bernard Badzioch, \emph{MTH 727 Lecture Notes}, Chapter 17, \emph{Hurewicz Theorem}。作者姓名和CC许可取自同目录 mth727\_notes\_header.tex。原文件17\_hurewicz\_theorem.tex，开头至 HUREWICZ ISOMORPHISM THM 证明结束（第1--256行）；Git blob SHA \texttt{594fdb6770ca9b03bb125949ce30a094f72ba176}；获取2026-09-28。
\par\url{https://github.com/bbadzioch/topology_lecture_notes/blob/master/MTH_727_notes/17_hurewicz_theorem.tex}
\par\url{https://github.com/bbadzioch/topology_lecture_notes/blob/master/MTH_727_notes/mth727_notes_header.tex}
\par\textbf{本文件只计一个问题（整理者概述）：}证明 $(n-1)$ 连通空间的首个可能非零同伦群由球面基本类诱导地同构于 $H_n$，$n\geq2$，并证明低阶同调消失。
\subsection*{作者命题与人类证明}

Hurewicz homomorphism is a map that connects homotopy and homology groups.
Recall that $H_n(S^n)\cong\Z$. We will denote by $\gamma_n$ a chosen generator of $H_n(S^n)$. Given an element $[\varphi:(S^n,s_0)\to(X,x_0)]\in\pi_n(X,x_0)$ consider the homomorphism $\varphi_*:H_*(S^n)\to H_n(X)$. This homomorphism depends only on the homotopy class of $\varphi$.
\begin{definition}
The \emph{Hurewicz homomorphism} is a function
\[h:\pi_n(X,x_0)\to H_n(X)\]
given by $h([\varphi])=\varphi_*(\gamma_n)$.
\end{definition}
\begin{proposition}[Naturality]
For any function $f:X\to Y$ the following diagram commutes:
\[
\xymatrix{\pi_n(X,x_0)\ar[r]^{f_*}\ar[d]_h&\pi_n(Y,f(x_0))\ar[d]^h\\H_n(X)\ar[r]_{f_*}&H_n(Y)}
\]
\end{proposition}
\begin{proof}
For $[\varphi]\in\pi_n(X,x_0)$ we have
\[hf_*([\varphi])=h([f\varphi])=(f\varphi)_*(\gamma_n)=f_*\varphi_*(\gamma_n)=f_*h([\varphi]).\]
\end{proof}
\begin{proposition}
The Hurewicz homorphism is a group homomorphism.
\end{proposition}
\begin{proof}
Let $\varphi,\psi:(S^n,s_0)\to(X,x_0)$ where $n\geq1$. Recall that the element $[\varphi]\cdot[\psi]\in\pi_n(X,x_0)$ is the homotopy class of the map
\[S^n\overset p\longrightarrow S^n\vee S^n\overset{\varphi\vee\psi}\longrightarrow X\]
where $p:S^n\to S^n\vee S^n$ is the pinch map. Let $r_1,r_2:S^n\vee S^n\to S^n$ be the retractions of $S^n\vee S^n$ onto the first and, respectively, the second copy of $S^n$. We have a commutative diagram
\[
\xymatrix@C=3em{H_n(S^n)\ar[r]^{p_*}\ar[dr]_{\id_*\oplus\id_*}&H_n(S^n\vee S^n)\ar[r]^{(\varphi\vee\psi)_*}\ar[d]_{r_{1*}\oplus r_{2*}}^{\cong}&H_n(X)\\&H_n(S^n)\oplus H_n(S^n)\ar[ur]_{\varphi_*+\psi_*}&}
\]
This gives:
\begin{align*}
h([\varphi]\cdot[\psi])&=((\varphi\vee\psi)p)_*(\gamma_n)\\
&=(\varphi_*+\psi_*)(\id_*\oplus\id_*)(\gamma_n)\\
&=\varphi_*(\gamma_n)+\psi_*(\gamma_n)=h([\varphi])+h([\psi]).
\end{align*}
\end{proof}
\begin{theorem}[Hurewicz isomorphism theorem]
Let $X$ be a path connected space such that for some $n\geq2$ we have $\pi_i(X)=0$ for $i<n$. Then $H_i(X)=0$ for $0<i<n$ and the Hurewicz homomorphism
\[h:\pi_n(X,x_0)\to H_n(X)\]
is an isomorphism.
\end{theorem}
\begin{proof}
Assume first that $X=S^n$. We have $H_i(S^n)=0$ for $0<i<n$. In degree $n$ the Hurewicz homomorphism is a map $h:\Z\cong\pi_n(S^n)\to H_n(S^n)\cong\Z$. The group $\pi_n(S^n)$ is generated by the homotopy class of the identity map $\id_{S^n}:S^n\to S^n$ (\textsf{PIN SN NOTE}). We have $h([\id_{S^n}])=\id_{S^n*}(\gamma_n)=\gamma_n$. Therefore $h$ maps a generator of $\pi_n(S^n)$ to a generator of $H^n(S^n)$, and so it is an isomorphism.

Next, assume that $X=\bigvee_{i\in I}S^n$. Again, in this case $H_i(\bigvee_{i\in I}S^n)=0$ for $0<i<n$. Also, the retraction maps $r_i:\bigvee_{i\in I}S^n\to S^n$ give a commutative diagram
\[
\xymatrix{\pi_n(\bigvee_{i\in I}S^n)\ar[r]^{\bigoplus r_{i*}}_{\cong}\ar[d]_h&\bigoplus_{i\in I}\pi_n(S^n)\ar[d]^{\bigoplus_{i\in I}h}_{\cong}\\H_n(\bigvee_{i\in I}S^n)\ar[r]_{\bigoplus r_{i*}}^{\cong}&\bigoplus_{i\in I}H_n(S^n)}
\]
The the map $\bigoplus_{i\in I}h$ is an isomorphism by the previous case, so the left vertical map $h$ is also an isomorphism.

For the next step, assume that $X$ is an arbitrary CW complex with $\pi_i(X)=0$ for $i<n$. By Proposition \textsf{N SKELETON FOR N CONNECTED PROP} we can assume that $X^{(n-1)}=*$, which gives $H_i(X)=0$ for $0<i<n$.

Let $j:X^{(n+1)}\hookrightarrow X$ be the inclusion of the $(n+1)$-skeleton of $X$. By the naturality proposition we have a commutative diagram
\[
\xymatrix{\pi_n(X^{(n+1)})\ar[r]^{j_*}_{\cong}\ar[d]_h&\pi_n(X)\ar[d]^h\\H_n(X^{(n+1)})\ar[r]_{j_*}^{\cong}&H_n(X)}
\]
The upper homomorphism $j_*$ is an isomorphism by Proposition \textsf{PIN FOR CW SKELETON PROP}, and the lower $j_*$ is an isomorphism by properties of homology groups. As a consequence, it is enough to show that $h:\pi_n(X^{(n+1)})\to H_n(X^{(n+1)})$ is an isomorphism.

Since $X^{(n-1)}=*$, it follows that $X^{(n)}=\bigvee_{i\in I}S^n$ and $X^{(n+1)}=X^{(n)}\cup\bigcup_{k\in K}e_k^{(n+1)}$ where $\{e_k^{(n+1)}\}_{k\in K}$ are $(n+1)$-cells of $X$. Let $\varphi_k:S^n\to X^{(n)}$ be the attaching map of the cell $e_k^{n+1}$, and let $i:X^{(n)}\hookrightarrow X^{(n+1)}$ denote the inclusion map. We have a commutative diagram
\[
\xymatrix@C=1.1em{
\pi_n(\bigvee_{k\in K}S^n)\ar[r]^{(\bigvee\varphi_k)_*}\ar[d]_h^{\cong}&\pi_n(X^{(n)})\ar[r]^{i_*}\ar[d]_h^{\cong}&\pi_n(X^{(n+1)})\ar[r]\ar[d]^h&0\ar[d]^{\cong}\\
H_n(\bigvee_{k\in K}S^n)\ar[r]_{(\bigvee\varphi_k)_*}&H_n(X^{(n)})\ar[r]_{i_*}&H_n(X^{(n+1)})\ar[r]&H_{n-1}(\bigvee_{k\in K}S^n)
}
\]
The upper row of this diagram is exact by Proposition \textsf{1CONN PIN XCUPE PROP}, and the lower row is exact by the long homology sequence associated to the map $\bigvee_{k\in K}\varphi_k$. By the Five Lemma we obtain that $h:\pi_n(X^{(n+1)})\to H_n(X^{(n+1)})$ is an isomorphism.

Finally, let $X$ be an arbitary space with $\pi_i(X)=0$ for $i<n$. Let $f:Y\to X$ be a CW approximation of $X$ (\textsf{CW APPROX DEF}). Using Theorem \textsf{WE HOMOLOGY ISO THM} and the previous case we get $H_i(X)\cong H_i(Y)=0$ for $0<i<n$.

By the naturality proposition we have a commutative diagram
\[
\xymatrix{\pi_n(Y)\ar[r]^{f_*}_{\cong}\ar[d]_h^{\cong}&\pi_n(X)\ar[d]^h\\H_n(Y)\ar[r]_{f_*}^{\cong}&H_n(X)}
\]
Since $f$ is a weak equivalence, the upper homomorphism $f_*$ is an isomorphism by definition, and the lower $f_*$ is an isomorphism by Theorem \textsf{WE HOMOLOGY ISO THM}. Also, since $Y$ is a CW complex the left vertical map is an isomorphism by the previous case. Therefore $h:\pi_n(X)\to H_n(X)$ is an isomorphism.
\end{proof}

\subsection*{整理说明与先修范围（非作者正文）}
完整保留 Hurewicz 映射定义、自然性、同态性质及主定理的球面、楔和、胞腔复形和一般空间四段证明；后面的逆Hurewicz定理、$\pi_1$阿贝尔化版本及练习没有取入。交换图用等价的XY-pic布局重排，箭头和节点保留；内部无法独立解析的引用改成原文可搜索的标签名。作者原文一次将 $H_n(S^n)$ 写作 $H^n(S^n)$、以及 homorphism/arbitary/the the 字样照录，不静默改数学内容。标准先修为胞腔逼近、同伦群对骨架及粘胞腔的正合描述、CW逼近及弱等价的同调不变性、同调长正合列和五引理。难点是选择 $(n+1)$ 骨架并把同伦粘接关系与同调边界关系放到同一个自然图中；这里没有把Hurewicz结论本身列作先修。源码由连接器取得，包内为整理转录摘录。
\subsection*{可修改的简短标注（非作者正文）}
$c$：空间首个非零同伦群的同调刻画

$a$：球面基本类；捏映射；CW逼近；骨架与粘胞腔；同伦及同调正合列；自然性；五引理

\texttt{cross\_theory\_status = yes}。将球面映射类送到基本类的像，用胞腔骨架把两类不变量中的生成元与关系逐一对接，最后由正合图返回同构。位置：Chapter 17 主定理的CW复形步骤和末段CW逼近。

全部标注为非穷尽、可修改初稿。
\subsection*{来源许可与修改记录}
原作者 Bernard Badzioch；仓库README与章节导言明确采用CC BY-NC-SA 4.0。本转录及整理部分采用同许可，保留作品、源文件及作者归属。
2026-09-28：ChatGPT 整理为独立 TeX，加入编号、来源定位、整理说明与初稿标注；原作者数学论证与整理者文字分开标示。
\end{document}
